%% OVERFORE-Formulario-T-10.tex — generado desde OVERFORE-Formulario-T-10.docx
%% Formato: clase overfore, modo formulario. El formato lo pone la clase; aquí sólo va el contenido.
\documentclass[formulario]{../../comun/overfore}
\numsubdoc{6}
\formulario{T-10}{Metodología para el desarrollo}
\informe{Informe Preparatorio N° 2}
\fecha{12 de octubre de 2026}

\begin{document}

\apartado{Metodología de Desarrollo Software}
En este apartado se presenta el resultado que persigue la metodología de desarrollo, la organización común de los equipos y la forma en que se autoriza el paso a producción.

\begin{tabla}{I{0.176}I{0.674}I{0.150}}{Aplicación de la metodología de desarrollo}{tab:aplicacion-metodologia-desarrollo}
  {\cab{Elemento} & \cab{Aplicación en SITCO Aranda} & \cab{Registro}}
  {Fuente: elaboración propia a partir del Subdocumento 6.}
Resultado esperado & Transformar un requerimiento en una versión probada y autorizada, conservando trazabilidad, aprobaciones y resultados. & Registro de versión \\
Organización común & Backlog, definición de terminado y flujo de entrega compartidos. Refinamiento con el cliente manteniendo aceptación y ventanas operacionales. & Backlog y criterios \\
Autorización & El comité de proyecto registra el paso a producción; el flujo automatizado ejecuta la instalación. La nube, el banco de pruebas y el nodo de faena conservan evidencias diferenciadas. & Aprobación y despliegue \\
\end{tabla}

\seccion{Enfoque de desarrollo}
Este apartado muestra el método de desarrollo adoptado, la organización de los equipos y la planificación de los sprints en cada etapa del contrato.

\begin{tabla}{I{0.190}I{0.628}I{0.182}}{Equipos y planificación del desarrollo}{tab:equipos-planificacion-desarrollo}
  {\cab{Elemento del enfoque} & \cab{Aplicación en SITCO Aranda} & \cab{Registro}}
  {Fuente: elaboración propia a partir del Subdocumento 6.}
Método & Scrum con tres equipos coordinados y sprints de dos semanas (Schwaber y Sutherland, 2020). Validación del incremento con ambos grupos 7x7. & Objetivo e incremento de sprint \\
Equipo de cadena del mineral & Dedicado a trazabilidad y conciliación. & Historias de esos dominios \\
Equipo de personas y cumplimiento & Dedicado a funciones de personas y cumplimiento. & Historias de esos dominios \\
Equipo de plataforma e integración & Prepara servicios e integraciones. El líder de integración responde por el integrador y por la coordinación con SAP, despacho, laboratorio y demás sistemas. & Dependencias y conectores \\
Etapa 1 & 19 sprints entre los meses 3 y 11; certificación en el mes 12. & Plan de construcción \\
Etapa 2 & Diez sprints entre los meses 13 y 17; certificación en el mes 18. Entre los meses 12 y 17, el paquete 1.4.3.7 corrige los defectos de la etapa 1, sin agregar requerimientos. & Plan y reserva de capacidad \\
Ciclo y operación & Preparación semanal, selección, construcción con controles, revisión y liberación aprobada. Desde el mes 21, Kanban de mantención conducido por el líder de operación (SRE). & Registros del flujo \\
\end{tabla}

\subseccion{Gestión de requerimientos}
Este apartado presenta cómo se descomponen, priorizan y trazan los requerimientos, y las condiciones que debe reunir una historia para entrar a un sprint.

\begin{tabla}{I{0.200}I{0.607}I{0.193}}{Preparación y trazabilidad de requerimientos}{tab:preparacion-trazabilidad-requerimientos}
  {\cab{Práctica} & \cab{Aplicación en SITCO Aranda} & \cab{Evidencia}}
  {Fuente: elaboración propia a partir del Subdocumento 6.}
Descomposición & La historia conserva el identificador del requerimiento. Épicas por módulo; requerimientos funcionales en historias; no funcionales en criterios de aceptación y definición de terminado. & Historia y criterios \\
Priorización & Etapa comprometida, dependencia técnica y valor para el usuario, en ese orden. El líder funcional y los referentes revisan semanalmente el backlog; el líder de datos acuerda con ellos la definición de cada dato y sus reglas de calidad. & Backlog ordenado \\
Cadena de evidencia & Requerimiento → tarea ClickUp → rama → solicitud de integración → ejecución → artefacto → despliegue. ClickUp y GitHub permiten recorrerla en ambos sentidos. & Identificadores vinculados \\
Botonera & Una prueba automática verifica el registro del cambio de material en dos segundos o menos, en el percentil 95. Se verifica en cada historia de la botonera durante todo el desarrollo. La aceptación funcional no sustituye la prueba de desempeño. & Resultado de la prueba \\
Cambios de aceptación & Si una historia modifica un compromiso de aceptación, se aplica el control formal de cambios. & Solicitud aprobada \\
\end{tabla}
\begin{tabla}{I{0.319}I{0.369}I{0.312}}{Condiciones de entrada de historias}{tab:condiciones-entrada-historias}
  {\cab{Dato de la historia} & \cab{Registro} & \cab{Verificación de entrada}}
  {\emph{Fuente: }elaboración propia a partir del Subdocumento 6.}
Requerimiento de origen & Identificador del catálogo & Vinculación conservada \\
Módulo y etapa & Épica y compromiso de entrega & Prioridad compatible \\
Criterio de aceptación & Resultado verificable & Prueba identificable \\
Dependencias & Sistemas y tareas relacionadas & Resueltas para iniciar \\
Datos de prueba & Datos disponibles para el escenario & Preparación comprobada \\
Decisión funcional & Acuerdo con referente del cliente & Refinamiento registrado \\
\end{tabla}

\subseccion{Arquitectura evolutiva}
Este apartado presenta las decisiones estructurales iniciales, los ensayos con que se verifica la arquitectura y las propiedades que se comprueban durante el desarrollo.

\begin{tabla}{I{0.176}I{0.644}I{0.180}}{Evolución y verificación de arquitectura}{tab:evolucion-verificacion-arquitectura}
  {\cab{Decisión o ensayo} & \cab{Aplicación en SITCO Aranda} & \cab{Evidencia}}
  {Fuente: elaboración propia a partir del Subdocumento 6.}
Acuerdos iniciales & El integrador, el nodo de faena y los servicios de identidad y auditoría se fijan en los primeros sprints. Los cambios transversales o tecnológicos requieren alternativas y aprobación del comité de arquitectura. & Decisión de arquitectura \\
Verificación continua & Las propiedades se comprueban mediante funciones de aptitud durante el desarrollo (Ford et al., 2017). & Resultados de verificación \\
Ensayo por versión & Cada versión del nodo de faena se ensaya en el banco de pruebas de la sala técnica. & Registro del banco \\
Aceptación física & Antes de la producción de cada etapa, corte controlado de 48 horas sobre el nodo de faena: continuidad local, conservación de registros y sincronización completa en máximo cuatro horas tras recuperar el enlace. & Informe de aceptación \\
Repetición operacional & Ensayo anual en detención programada de planta. Se mantienen separados los registros de autonomía física y los de desempeño de preproducción. & Informe anual y resultados separados \\
Alcance de la prueba & El corte real no se sustituye por simulación en nube ni se repite obligatoriamente antes de cada liberación ordinaria. & Plan de pruebas \\
\end{tabla}
\begin{tabla}{I{0.241}I{0.316}I{0.251}I{0.192}}{Verificaciones de arquitectura}{tab:verificaciones-arquitectura}
  {\cab{Propiedad} & \cab{Lugar de verificación} & \cab{Condición comprobada} & \cab{Evidencia}}
  {Fuente: elaboración propia a partir del Subdocumento 6.}
Autonomía local & Banco de pruebas de la sala técnica y nodo de faena & Continuidad sin enlace & Resultados del ensayo \\
Registro en terreno & Aplicación y dispositivo & Tiempo comprometido & Prueba de desempeño \\
Compatibilidad de integración & Simuladores de contratos & Interfaz compatible & Resultado de contrato \\
Separación del control & Componentes e interfaces & Ausencia de escritura a red de control & Verificación \\
\end{tabla}

\subseccion{Refactorización y deuda técnica}
Este apartado presenta el registro, la reserva de capacidad y los umbrales con que se controla la deuda técnica.

\begin{tabla}{I{0.144}I{0.687}I{0.169}}{Control de deuda técnica}{tab:control-deuda-tecnica}
  {\cab{Control} & \cab{Regla aplicada} & \cab{Evidencia}}
  {Fuente: elaboración propia a partir del Subdocumento 6.}
Registro y revisión & El autor deja el código más claro y mantenible; los pares comprueban el cambio. La deuda registra causa, impacto, esfuerzo y módulo afectado. & Registro de deuda y revisión \\
Reserva & 15 \% de la capacidad de cada sprint para deuda técnica. La deuda de seguridad o desempeño se resuelve en el sprint siguiente a su detección. & Plan del sprint \\
Revisión y umbral & El comité de arquitectura revisa mensualmente. Un índice de deuda mayor que 5 \% del esfuerzo del código nuevo de un módulo bloquea el cambio en las puertas 1 y 2. Exige dedicar el sprint siguiente a reducirlo. & Análisis estático y acciones \\
Bloqueo de seguridad & No se promueve con hallazgos críticos o altos pendientes, aunque otras mejoras sigan planificadas. & Controles aprobados \\
\end{tabla}

\subseccion{Tiempo de salida al mercado}
Este apartado presenta las condiciones que debe cumplir un incremento para pasar de la revisión a producción, considerando las restricciones de la operación en faena.

\begin{tabla}{I{0.176}I{0.664}I{0.160}}{Revisión y condiciones de liberación}{tab:revision-condiciones-liberacion}
  {\cab{Condición} & \cab{Regla aplicada} & \cab{Evidencia}}
  {Fuente: elaboración propia a partir del Subdocumento 6.}
Incremento & Despliegue automático del incremento de cada sprint en preproducción para revisión. & Resultado de revisión \\
Usuarios de faena & Desde el mes 5, pruebas con usuarios de faena, con guantes, sol y zonas sin señal; desde el mes 8, con equipos reales. Resultados al backlog antes de la marcha blanca. & Validaciones de uso \\
Ventana permitida & No coincidir con cierre contable, 48 horas previas a embarque ni detención de planta. Cambios en componentes de faena entre 03:00 y 05:00. & Autorización de ventana \\
Ensayo anual & La detención programada usada para el ensayo de autonomía es una actividad específica de prueba, separada del despliegue ordinario. & Plan de ensayo \\
Autorización separada & Revisar un incremento no autoriza producción: artefacto técnicamente disponible y aprobación operacional son condiciones diferentes. & Aprobación de liberación \\
\end{tabla}

\seccion{Prácticas de DevSecOps, integración y entrega continuas}
Este apartado presenta las reglas de seguridad que se aplican a todo el flujo de entrega y que condicionan las prácticas descritas en los apartados siguientes.

\begin{tabla}{I{0.186}I{0.578}I{0.236}}{Reglas de seguridad del flujo}{tab:reglas-seguridad-flujo}
  {\cab{Regla} & \cab{Aplicación en el flujo} & \cab{Evidencia}}
  {Fuente: elaboración propia a partir del Subdocumento 6.}
Marco de trabajo & Seguridad incorporada al cambio mediante controles y decisiones registradas, conforme al SSDF (NIST, 2022). & Resultados vinculados a versión \\
Bloqueo por severidad & Los hallazgos críticos o altos pendientes impiden integrar o promover. No se omiten controles bloqueantes. & Informes aprobados \\
Otros bloqueos & Credenciales expuestas, pruebas fallidas, licencias incompatibles y cobertura insuficiente detienen el flujo. & Registro de rechazo y corrección \\
\end{tabla}

\subseccion{Repositorio y flujo de trabajo}
Este apartado presenta las reglas de titularidad, ramas, revisión y versionado del código, y la organización de los repositorios.

\begin{tabla}{I{0.164}I{0.660}I{0.177}}{Reglas de repositorio y revisión}{tab:reglas-repositorio-revision}
  {\cab{Regla} & \cab{Aplicación en SITCO Aranda} & \cab{Evidencia}}
  {Fuente: elaboración propia a partir del Subdocumento 6.}
Titularidad & Organización GitHub Enterprise Cloud de propiedad del cliente, con su código e historial; Overfore participa como miembro. & Organización y permisos \\
Ramas & Rama corta de menos de dos días con identificador ClickUp; commits firmados y cambios a la rama principal solo por solicitudes aprobadas (Humble y Farley, 2010). & Rama y solicitud \\
Revisión reforzada & CODEOWNERS y reglas de protección exigen dos revisores para cambios en el nodo de faena, la seguridad o el cálculo crítico. & Aprobaciones registradas \\
Versionado y acceso & Etiquetas de versionado semántico para producción y correcciones urgentes; identidades federadas con Azure, sin credenciales permanentes en el flujo. & Versión e identidad \\
Corrección & Un control fallido se corrige en la rama y se reejecuta. Solo se avanza con aprobaciones y resultados válidos. & Tarea, revisiones y ejecuciones \\
\end{tabla}
\begin{tabla}{I{0.385}I{0.332}I{0.283}}{Organización de repositorios}{tab:organizacion-repositorios}
  {\cab{Repositorio} & \cab{Contenido administrado} & \cab{Uso en el flujo}}
  {Fuente: elaboración propia a partir del Subdocumento 6.}
sitco-aplicacion & Nueve módulos, integrador y portales & Construcción de servicios \\
sitco-terreno & SITCO Terreno y botonera & Construcción de software de terreno \\
sitco-infraestructura & OpenTofu de nube, nodo de faena y borde & Plan de infraestructura \\
sitco-despliegue & Configuración y versiones de ambientes & Promoción y actualización \\
\end{tabla}

\subseccion{Prácticas de integración continua}
Este apartado presenta la secuencia de la integración continua, el tratamiento de sus resultados y los controles que un cambio debe superar para avanzar.

\begin{tabla}{I{0.213}I{0.625}I{0.162}}{Secuencia y tratamiento de resultados de integración}{tab:secuencia-tratamiento-resultados-integracion}
  {\cab{Práctica de integración} & \cab{Aplicación en SITCO Aranda} & \cab{Evidencia}}
  {Fuente: elaboración propia a partir del Subdocumento 6.}
Ejecución y amenazas & Cada solicitud y cada cambio a la rama principal ejecutan CI. Antes de construir un cambio arquitectónico se actualiza el modelo STRIDE como preparación de la historia (Shostack, 2014). & Modelo y ejecución \\
Secretos y cobertura & Escaneo de envíos e historial; una credencial expuesta detiene el flujo y avisa. Las pruebas unitarias verifican la lógica de negocio; una cobertura menor de 70 \% no se suple manualmente. & Escaneo y pruebas \\
Análisis de código y componentes & Se verifican calidad, seguridad, deuda, vulnerabilidades, licencias y lista de materiales de cada versión. & Informes y SBOM \\
Infraestructura & OpenTofu se analiza por exposición de almacenamiento, puertos, cifrado, políticas y etiquetas. Una etiqueta faltante incumple la política; no es automáticamente un hallazgo crítico de seguridad. & Análisis del plan \\
Contratos & El líder de integración mantiene contratos y simuladores para no depender de SAP, despacho o laboratorio reales en las pruebas diarias. & Resultados por conector \\
Artefacto único & Construir una vez, escanear, aprobar, firmar con procedencia y publicar. La revisión comprueba el cambio y el plan de infraestructura cuando corresponde. No se recompila entre ambientes. & Imagen firmada y registro \\
Rechazo & Se detiene y notifica con resultado vinculado a la tarea; el responsable corrige y vuelve al flujo. Aprobar exige resultados y revisiones, no solo una imagen publicada. & Reejecución y aprobaciones \\
\end{tabla}
\begin{tabla}{I{0.246}I{0.198}I{0.358}I{0.198}}{Controles de integración continua}{tab:controles-integracion-continua}
  {\cab{Control} & \cab{Herramienta} & \cab{Condición para avanzar} & \cab{Evidencia}}
  {Fuente: elaboración propia a partir del Subdocumento 6.}
Preparación del trabajo & GitHub Actions & Código y ejecución preparados & Ejecución \\
Secretos & GitHub Secret Protection & Sin credenciales expuestas & Resultado del escaneo \\
Compilación y pruebas & GitHub Actions & Pruebas aprobadas y cobertura de negocio ≥ 70 \% & Informe de pruebas \\
Análisis estático & SonarQube Cloud & Calidad aprobada y sin hallazgos críticos ni altos & Informe de análisis \\
Dependencias y licencias & Dependabot y Trivy & Sin vulnerabilidades críticas o altas y licencias compatibles & Informe y SBOM \\
Infraestructura como código & Checkov & Plan conforme y sin hallazgos críticos ni altos & Análisis de OpenTofu \\
Contratos & Pact y WireMock & Interfaz compatible & Resultados por conector \\
Imagen & Docker y Trivy & Escaneo aprobado sin hallazgos críticos ni altos & Identificador e informe \\
Firma y procedencia & Cosign & Imagen aprobada antes de firmar & Firma y procedencia \\
Revisión por pares & CODEOWNERS y GitHub & Aprobaciones requeridas & Solicitud aprobada \\
Publicación & Azure Container Registry & Artefacto firmado aprobado & Registro de artefacto \\
\end{tabla}

\subseccion{Prácticas de entrega continua}
Este apartado presenta las condiciones de promoción entre ambientes, los mecanismos de recuperación y los controles de la entrega continua en cada destino.

\begin{tabla}{I{0.213}I{0.620}I{0.167}}{Condiciones de promoción y recuperación}{tab:condiciones-promocion-recuperacion}
  {\cab{Práctica de entrega} & \cab{Aplicación en SITCO Aranda} & \cab{Evidencia}}
  {Fuente: elaboración propia a partir del Subdocumento 6.}
Promoción & Mismo artefacto firmado en todos los ambientes; autorización humana en ambiente protegido y ejecución automatizada. Firma comprobada antes de instalar. & Versión y autorización \\
Infraestructura y secretos & OpenTofu se compara cada noche con el código y alerta desviaciones. Configuración separada por ambiente; secretos en Key Vault con rotación y auditoría. & Plan, alerta y accesos \\
Uso de ambientes & Ambientes efímeros vinculados a ramas; no productivos reducidos fuera de horario. QA prueba integración, extremo a extremo y análisis dinámico; preproducción verifica desempeño, resiliencia y despliegue azul-verde. & Resultados de ambiente \\
Ensayos físicos & El banco de pruebas verifica cada versión. El nodo de faena se prueba 48 horas antes del paso a producción de cada etapa y anualmente en operación; sincronización máxima de cuatro horas. & Registros físicos separados \\
Migración & El líder de datos aprueba migraciones versionadas, automáticas y reversibles; ambas versiones de la aplicación conviven durante el cambio de esquema. & Aprobación y reversión \\
Aprobación productiva & Pruebas de aceptación, pruebas de seguridad, intrusión independiente, plan de reversión probado, ventana permitida y usuarios capacitados. & Expediente de liberación \\
Destinos & Nube con azul-verde. Flux, mediante GitOps de Azure Arc, descarga en faena cuando hay enlace y aplica en la ventana conservando la versión anterior. Botonera por anillos: una pala y después las cinco. App SITCO Terreno por anillos: un área y un turno, el resto del área y luego las demás áreas. & Resultado por destino \\
Recuperación & Una falla previa detiene la promoción; una falla instalada activa la reversión. El líder de operación (SRE) vigila indicadores y degradación posterior. & Monitoreo y reversión \\
Remediación operacional & Críticos en siete días corridos, altos en quince y medios en treinta. Estos plazos no autorizan versiones nuevas con hallazgos críticos o altos pendientes. & Hallazgos y corrección \\
\end{tabla}
\begin{tabla}{I{0.247}I{0.198}I{0.356}I{0.198}}{Controles de entrega continua}{tab:controles-entrega-continua}
  {\cab{Actividad} & \cab{Destino} & \cab{Condición para avanzar} & \cab{Evidencia}}
  {Fuente: elaboración propia a partir del Subdocumento 6.}
Verificar artefacto & Todos los ambientes & Firma y procedencia válidas & Verificación \\
Aplicar infraestructura & Nube y faena & Plan OpenTofu aprobado & Plan y ejecución \\
Pruebas funcionales y de seguridad & QA & Integración, extremo a extremo y seguridad aprobados & Resultados \\
Desempeño y resiliencia & Preproducción & Tiempos, fallas y azul-verde verificados & Informe \\
Ensayo de versión & Banco de pruebas en sala técnica & Comportamiento del nodo de faena verificado & Registro del ensayo \\
Aceptación de autonomía & Nodo de faena & 48 h y sincronización ≤ 4 h & Informe de aceptación \\
Migración de esquema & Ambiente de instalación & Migración reversible aprobada por el líder de datos & Versión y ejecución \\
Autorizar producción & Ambiente protegido & Comité de proyecto y condiciones de liberación aprobados & Aprobación registrada \\
Desplegar en nube & Producción & Azul-verde con artefacto aprobado & Despliegue \\
Desplegar en nodo de faena & Faena & Enlace disponible y ventana de 03:00 a 05:00 & Actualización por Flux \\
Desplegar en terreno & Palas y tabletas & Primer anillo verificado antes del siguiente & Resultado por anillo \\
Verificar operación & Nube y faena & Indicadores dentro de lo comprometido & Monitor y resultado \\
\end{tabla}

\subseccion{Herramientas de integración continua}
Este apartado presenta las herramientas de la integración continua y la evidencia que cada una deja en el expediente de versión.

\begin{tabla}{I{0.298}I{0.552}I{0.150}}{Herramientas y evidencia de integración}{tab:herramientas-evidencia-integracion}
  {\cab{Herramienta} & \cab{Uso en el flujo} & \cab{Evidencia generada}}
  {Fuente: elaboración propia a partir del Subdocumento 6.}
GitHub Actions & Ejecución del flujo vinculada a solicitud y tarea. & Ejecución \\
Secret Protection, SonarQube Cloud, Dependabot, Trivy, Checkov y OWASP ZAP & Resultados de seguridad, calidad, dependencias, infraestructura y análisis dinámico con OWASP ZAP. & Informes y bloqueos \\
Pact y WireMock & Compatibilidad de contratos de integración. & Resultados por conector \\
Docker y Cosign & Construcción de imagen y conservación de firma y procedencia. & Artefacto firmado \\
Azure Container Registry & Publicación del artefacto firmado y aprobado. & Registro de artefacto \\
CODEOWNERS y GitHub & Revisiones y aprobaciones exigidas. & Solicitud aprobada \\
Integración de ClickUp con GitHub & Localizar tarea, solicitud, informes, identificador de imagen y registro en el Container Registry sobre el mismo cambio y artefacto. & Expediente de versión \\
\end{tabla}

\subseccion{Herramientas de entrega continua}
Este apartado presenta las herramientas de la entrega continua y el registro que generan en cada liberación.

\begin{tabla}{I{0.287}I{0.522}I{0.191}}{Herramientas y evidencia de entrega}{tab:herramientas-evidencia-entrega}
  {\cab{Herramienta o registro} & \cab{Uso en el flujo} & \cab{Evidencia generada}}
  {Fuente: elaboración propia a partir del Subdocumento 6.}
GitHub Actions y Cosign & Ambientes protegidos, aprobación y verificación del artefacto. & Autorización y verificación \\
OpenTofu, Key Vault e identidades federadas & Aplicación de infraestructura y administración de acceso. & Plan, ejecución y acceso \\
Datafaker & Preparación de datos sintéticos. & Datos de prueba \\
Playwright, OWASP ZAP, k6 y Azure Chaos Studio & Pruebas de QA y preproducción. & Resultados de pruebas \\
Flux (GitOps de Azure Arc) y Azure Monitor & Versiones al nodo de faena y comprobación operacional. & Actualización y monitoreo \\
Registro de liberación & Versión, ambiente, aprobación y resultado; identifica la versión anterior para comprobar la reversión. & Registro de instalación \\
\end{tabla}

\subseccion{Automatización de pruebas y métricas de entrega}
Este apartado presenta los niveles de prueba automatizada, la forma de medir el flujo de entrega y las metas para la implementación y la operación.

\begin{tabla}{I{0.176}I{0.679}I{0.145}}{Pruebas y medición del flujo}{tab:pruebas-medicion-flujo}
  {\cab{Nivel de prueba o medición} & \cab{Aplicación en SITCO Aranda} & \cab{Evidencia}}
  {Fuente: elaboración propia a partir del Subdocumento 6.}
Integración continua & Pruebas unitarias, de contrato y de análisis dinámico de seguridad, con resultado asociado a la ejecución y la versión. & Informes de CI \\
Entrega continua & Pruebas de integración, extremo a extremo, desempeño y resiliencia en sus ambientes. & Informes de CD \\
Faena & Ensayos físicos con registros propios, separados de preproducción. & Resultados de faena \\
Métricas & Medición desde el inicio conforme a Forsgren et al. (2018), con marcas de cambio y despliegue, resultados y recuperación. & Registro de métricas \\
Desviaciones & Se analizan las evidencias para identificar mejoras. La frecuencia de despliegue respeta las ventanas y no omite pruebas ni aprobación. & Acción de mejora \\
\end{tabla}
\begin{tabla}{I{0.312}I{0.357}I{0.331}}{Metas del flujo de entrega}{tab:metas-flujo-entrega}
  {\cab{Métrica} & \cab{Meta durante la implementación} & \cab{Meta durante la operación}}
  {Fuente: elaboración propia a partir del Subdocumento 6.}
Frecuencia de despliegue & Diaria en QA y al menos una por sprint en preproducción & Al menos una cada dos semanas en producción \\
Tiempo desde el cambio hasta el despliegue & Menos de un día hasta QA & Cinco días hábiles o menos hasta producción \\
Tasa de cambios fallidos & 5 \% o menos en preproducción & 5 \% o menos en producción \\
Tiempo de restauración & Menos de un día en ambientes de prueba & Cuatro horas o menos ante incidente crítico \\
\end{tabla}

\seccion{Ceremonias, artefactos, cadencias y mecanismos de decisión}
Este apartado presenta los artefactos del desarrollo, la autoridad de cada instancia de decisión y el calendario de ceremonias con sus participantes.

\begin{tabla}{I{0.213}I{0.627}I{0.160}}{Artefactos y autoridad de decisión}{tab:artefactos-autoridad-decision}
  {\cab{Artefacto o decisión} & \cab{Regla aplicada} & \cab{Registro}}
  {Fuente: elaboración propia a partir del Subdocumento 6.}
Responsabilidad del producto & El líder funcional, con los referentes del cliente. Revisión del incremento y aceptación de historias según criterios; el comité de proyecto autoriza producción. & Criterios y decisiones \\
Validación 7x7 & La validación semanal cubre ambos grupos en el sprint; se conservan resultados y correcciones de cada grupo. & Registro de validación \\
Artefactos & Backlog del producto y del sprint, incremento, definiciones de listo y de terminado, registro de decisiones de arquitectura y registro de deuda técnica. & Registros actualizados \\
Definición de listo & Requerimiento de origen, criterio de aceptación verificable, dependencias resueltas y datos de prueba disponibles. & Historia preparada \\
Definición de terminado & Revisión por pares, pruebas y cobertura aprobadas, sin hallazgos de seguridad críticos ni altos pendientes, documentación actualizada, despliegue automático en preproducción y trazabilidad al requerimiento de origen. El referente acepta o rechaza contra el criterio. & Historia y evidencias \\
Autoridad de decisión & El equipo y el líder funcional seleccionan el sprint; el referente acepta la función; el equipo implementa dentro de la arquitectura vigente; el comité de arquitectura aprueba cambios transversales; el comité de proyecto autoriza producción. & Decisión por instancia \\
Mantención & Desde el mes 21, Kanban conducido por el líder de operación (SRE). La evolutiva se ordena por la bolsa anual priorizada por el cliente, con límites de trabajo en curso; la correctiva sigue la severidad, conservando controles y definición de terminado. & Registro de mantención \\
\end{tabla}
La Tabla \ref{tab:plantilla-historia-usuario} es la plantilla de cada historia del backlog. Una historia entra a un sprint sólo con todos sus campos completos, que es la condición de listo.

\begin{tabla}{I{0.250}I{0.500}I{0.250}}{Plantilla de historia de usuario}{tab:plantilla-historia-usuario}
  {\cab{Campo} & \cab{Contenido} & \cab{Quién lo completa}}
  {Fuente: elaboración propia a partir del Subdocumento 6.}
Identificador y origen & Código de la historia y requerimiento de origen (RF o RNF del Formulario T-12), que se conserva hasta la prueba & Líder Funcional \\
Módulo y equipo & Módulo de SITCO y equipo Scrum que la construye & Líder de Desarrollo \\
Necesidad & Quién, qué necesita y para qué, en el lenguaje del área usuaria & Líder Funcional con el área \\
Criterios de aceptación & Casos verificables, incluidos los RNF como condición (por ejemplo, registro en 2 s al percentil 95) & Líder Funcional y Calidad \\
Dependencias & Historias, conectores o datos de otro equipo o del cliente que deben estar antes & Líder de Desarrollo \\
Estimación & Puntos del equipo y riesgo técnico declarado & Equipo \\
Evidencia de terminado & Pruebas automatizadas, revisión de código, análisis de seguridad y demostración aceptada & Equipo y Calidad \\
\end{tabla}

\begin{tabla}{I{0.248}I{0.256}I{0.259}I{0.237}}{Ceremonias y cadencias}{tab:ceremonias-cadencias}
  {\cab{Ceremonia} & \cab{Cadencia} & \cab{Participación} & \cab{Resultado}}
  {Fuente: elaboración propia a partir del Subdocumento 6.}
Planificación del sprint & Inicio del sprint; hasta dos horas & Equipo y líder funcional & Objetivo e historias \\
Reunión diaria & Diaria; quince minutos & Equipo & Coordinación e impedimentos \\
Refinamiento & Semanal; hasta noventa minutos & Equipo y referentes & Historias listas \\
Validación en terreno & Semanal & Usuarios del grupo en turno & Resultados de uso \\
Coordinación entre equipos & Dos veces por semana & Equipos & Dependencias resueltas \\
Revisión del sprint & Quincenal en preproducción & Referentes y usuarios clave & Aceptación o rechazo \\
Retrospectiva & Cierre del sprint & Equipo & Mejora para el sprint siguiente \\
Planificación de liberación & Antes de cada hito de entrega & Equipos y gobierno del proyecto & Entrega ordenada \\
\end{tabla}

\seccionsinnumero{Referencias}
\begin{referencias}
\item Ford, N., Parsons, R., y Kua, P. (2017). \emph{Building evolutionary architectures: Support constant change}. O'Reilly Media.
\item Forsgren, N., Humble, J., y Kim, G. (2018). \emph{Accelerate: The science of lean software and DevOps}. IT Revolution Press.
\item Humble, J., y Farley, D. (2010). \emph{Continuous delivery: Reliable software releases through build, test, and deployment automation}. Addison-Wesley.
\item National Institute of Standards and Technology. (2022). \emph{Secure Software Development Framework (SSDF) version 1.1: Recommendations for mitigating the risk of software vulnerabilities} (NIST SP 800-218). https://doi.org/10.6028/NIST.SP.800-218
\item Schwaber, K., y Sutherland, J. (2020). \emph{La Guía Scrum}. https://scrumguides.org/
\item Shostack, A. (2014). \emph{Threat modeling: Designing for security}. Wiley.
\end{referencias}
\end{document}
